\documentclass[11pt]{article}
\usepackage{mathblatt}
% gesetzt von werkzeuge/setzer.py (Sorte selbst) aus dem Lernweg MZE
% und den Bankzeilen; Muster bau/proben/2026-10-09/hypotenuse-muster.
\usepackage{enumitem}
\usepackage{adjustbox,varwidth}
\newgeometry{a4paper,margin=18mm,top=15mm,bottom=20mm,footskip=9mm}
\pagestyle{fancy}\fancyhf{}\renewcommand{\headrulewidth}{0pt}
\fancyfoot[L]{\footnotesize\color{mbgrau}MZE \textperiodcentered{} Lösungen \textperiodcentered{} Termwerte berechnen}
\fancyfoot[R]{\footnotesize\color{mbgrau}\thepage}
\setlength{\parskip}{3pt}
\renewcommand{\kreuz}[1]{\mbox{$\square$\ #1}\quad}
\newcommand{\kreuzl}[1]{\par\hangindent1.4em\hangafter1\noindent$\square$\ #1}
\newcommand{\aufg}[3]{\par\noindent\makebox[0pt][r]{#3}\begin{minipage}[t]{\linewidth}#1\end{minipage}\par}
\newcommand{\aufgg}[4]{\par\noindent\makebox[0pt][r]{#4}\begin{minipage}[t]{0.6\linewidth}\vspace{0pt}#1\end{minipage}\hfill
  \begin{minipage}[t]{0.37\linewidth}\vspace{0pt}\centering #2\end{minipage}\par}
\newcommand{\nr}[1]{\makebox[7mm][l]{\textbf{#1}}}
\newcommand{\lsg}[3]{\par\noindent\begin{minipage}[t]{9mm}\textbf{#1}\end{minipage}%
  \begin{minipage}[t]{52mm}\raggedright\bfseries\boldmath #2\end{minipage}\hspace{3mm}%
  \begin{minipage}[t]{\dimexpr\linewidth-64mm\relax}\raggedright\small #3\end{minipage}\par\vspace{4pt}}
\newlength{\kr}
\makeatletter
\newcommand{\karomerk}[2]{\protected@write\@auxout{}{\string\karoseite{#1}{#2}{\thepage}}}
\newcommand{\karoseite}[3]{\expandafter\xdef\csname karo@#1@#2\endcsname{#3}}
\newcommand{\karohinweis}[3]{\karomerk{h}{#1}\ifcsname karo@k@#1\endcsname
  \ifnum\csname karo@k@#1\endcsname=0\csname karo@h@#1\endcsname\relax#2\else#3\fi
  \else#3\fi}
\makeatother
\begin{document}

\noindent{\Large\bfseries Lösungen}\hspace{1em}{\color{mbgrau}Termwerte berechnen}\par\smallskip
\noindent{\small Links das Ergebnis, rechts der Weg in kurzen Schritten. Stimmt dein Ergebnis nicht, such den ersten Schritt, der anders ist.}\par
\par\Needspace{25mm}\vspace{6pt}\noindent\colorbox{black!12}{\makebox[7mm]{\bfseries\strut A}}\hspace{2mm}{\bfseries Zahl einsetzen und ausrechnen}\par\vspace{4pt}
\lsg{1.}{$6;\ 11;\ 46$}{a) $5 \cdot 2 - 4 = 10 - 4 = 6$; b) $5 \cdot 3 - 4 = 15 - 4 = 11$; c) $5 \cdot 10 - 4 = 50 - 4 = 46$}
\lsg{2.}{$13;\ 20;\ 11;\ 8$}{a) $7 + 2 \cdot 3 = 7 + 6 = 13$; b) $4 \cdot (3 + 2) = 4 \cdot 5 = 20$; c) $20 - 3 \cdot 3 = 20 - 9 = 11$; d) $2 \cdot 3 + 3 - 1 = 6 + 3 - 1 = 8$}
\lsg{3.}{$1;\ 4;\ 7;\ 10;\ 13$}{$3 \cdot 1 - 2 = 1$; $3 \cdot 2 - 2 = 4$; $3 \cdot 3 - 2 = 7$; $3 \cdot 4 - 2 = 10$; $3 \cdot 5 - 2 = 13$\par\smallskip \adjustbox{max width=\linewidth,max height=4cm}{\begin{varwidth}{50cm}\wertetabelle[1,4,7,10,13]{x}{3x - 2}{1,2,3,4,5}\end{varwidth}}}
\lsg{4.}{nein, $22$\,€; es fehlen $2$\,€}{$4 + 3 \cdot 6 = 4 + 18 = 22$\,€. Das Geld reicht nicht, es fehlen $22 - 20 = 2$\,€.}
\par\Needspace{25mm}\vspace{6pt}\noindent\colorbox{black!12}{\makebox[7mm]{\bfseries\strut B}}\hspace{2mm}{\bfseries Negative Zahlen einsetzen}\par\vspace{4pt}
\lsg{1.}{$4;\ -5;\ 7$}{a) $3 \cdot (-2) + 10 = -6 + 10 = 4$; b) $3 \cdot (-5) + 10 = -15 + 10 = -5$; c) $3 \cdot (-1) + 10 = -3 + 10 = 7$}
\lsg{2.}{$11;\ 6;\ 17;\ -9$}{a) $8 - (-3) = 8 + 3 = 11$; b) $-2 \cdot (-3) = 6$; c) $5 - 4 \cdot (-3) = 5 - (-12) = 5 + 12 = 17$; d) $-3 - 6 = -9$}
\lsg{3.}{$2;\ 18$}{a) $6 \cdot (-2 - 1) + 20 = 6 \cdot (-3) + 20 = -18 + 20 = 2$; b) $10 - 2 \cdot (-7 + 3) = 10 - 2 \cdot (-4) = 10 - (-8) = 18$}
\lsg{4.}{$3 - 4x$ (Wert $11$), um $3$ größer}{$3 - 4 \cdot (-2) = 3 + 8 = 11$; $2 \cdot (-2 + 6) = 2 \cdot 4 = 8$. $3 - 4x$ hat den größeren Wert, um $11 - 8 = 3$.}
\par\Needspace{25mm}\vspace{6pt}\noindent\colorbox{black!12}{\makebox[7mm]{\bfseries\strut C}}\hspace{2mm}{\bfseries Quadrate}\par\vspace{4pt}
\lsg{1.}{$25;\ 25;\ 1;\ 100$}{a) $5^2 = 25$; b) $(-5)^2 = (-5) \cdot (-5) = 25$; c) $(-1)^2 = 1$; d) $(-10)^2 = 100$}
\lsg{2.}{$16;\ 8;\ 6$}{a) $4 \cdot (-2)^2 = 4 \cdot 4 = 16$; b) $(-2)^2 + 4 = 4 + 4 = 8$; c) $(-2)^2 - (-2) = 4 + 2 = 6$}
\lsg{3.}{$1;\ -3$}{a) $10 - (-3)^2 = 10 - 9 = 1$; b) $2 \cdot (-1)^2 + 5 \cdot (-1) = 2 \cdot 1 + (-5) = 2 - 5 = -3$}
\lsg{4.}{$8;\ 3;\ 0;\ -1;\ 0;\ 3$; negativ nur für $x = 1$}{$(-2)^2 - 2 \cdot (-2) = 4 + 4 = 8$; $(-1)^2 - 2 \cdot (-1) = 1 + 2 = 3$; $0^2 - 2 \cdot 0 = 0$; $1^2 - 2 \cdot 1 = 1 - 2 = -1$; $2^2 - 2 \cdot 2 = 0$; $3^2 - 2 \cdot 3 = 9 - 6 = 3$. Negativ nur für $x = 1$.\par\smallskip \adjustbox{max width=\linewidth,max height=4cm}{\begin{varwidth}{50cm}\wertetabelle[8,3,0,-1,0,3]{x}{x^2 - 2x}{-2,-1,0,1,2,3}\end{varwidth}}}
\par\Needspace{25mm}\vspace{6pt}\noindent\colorbox{black!12}{\makebox[7mm]{\bfseries\strut D}}\hspace{2mm}{\bfseries Zwei Variablen und Bruchstrich}\par\vspace{4pt}
\lsg{1.}{$1;\ 11;\ -12;\ -7$}{a) $4 + (-3) = 1$; b) $2 \cdot 4 - (-3) = 8 + 3 = 11$; c) $4 \cdot (-3) = -12$; d) $-3 - 4 = -7$}
\lsg{2.}{$-2$}{$\dfrac{1 - (-7)}{-4} = 8 : (-4) = -2$}
\lsg{3.}{$-2{,}5$}{$\dfrac{6 + (-1)}{-2} = 5 : (-2) = -2{,}5$}
\lsg{4.}{$b - a$ (Wert $6$)}{$4 - (-2) = 6$. Die anderen: $a \cdot b = (-2) \cdot 4 = -8$; $\dfrac{b}{a} = 4 : (-2) = -2$; $a - b = -2 - 4 = -6$.}
\par\Needspace{25mm}\vspace{6pt}\noindent\colorbox{black!12}{\makebox[7mm]{\bfseries\strut T}}\hspace{2mm}{\bfseries Probetest}\par\vspace{4pt}
\lsg{1.}{$14$}{$6 - 2 \cdot (-4) = 6 - (-8) = 6 + 8 = 14$}
\lsg{2.}{$-20$}{$4 \cdot (-3 - 2) = 4 \cdot (-5) = -20$}
\lsg{3.}{$15$}{$2 \cdot (-3)^2 + (-3) = 2 \cdot 9 - 3 = 18 - 3 = 15$}
\lsg{4.}{$1{,}5$}{$\dfrac{-1 - 5}{-4} = (-6) : (-4) = 1{,}5$}
\lsg{5.}{$14$}{$(-3)^2 - (-5) = 9 + 5 = 14$}
\lsg{6.}{a) $18$\,m, ja \quad b) $40$\,m, nein}{a) $\dfrac{30}{10} = 3$; $3 \cdot 3 + 3^2 = 9 + 9 = 18$\,m; $18 < 20$: ja, es hält 2\,m vor dem Kind. b) $\dfrac{50}{10} = 5$; $3 \cdot 5 + 5^2 = 15 + 25 = 40$\,m; $40 > 20$: nein, der Weg ist 20\,m zu lang.}
\end{document}
